\subsection{Exponential strict indices with positive initial multiplicity}

\begin{proposition}[P-F62: odd dilations preserve the dyadic obstruction]
\label{prop:positivestartdilation}
Use the actual cubic stretching functions $f_n$ of P-F57.
Fix real numbers $\epsilon\ge0$ and $c>0$ satisfying
$\epsilon+c<1/4$. For all sufficiently large $n$, uniformly
over integers
\[
           1\le t\le2^{\epsilon n},\qquad \gcd(t,6)=1,
                                                               \tag{P-F62a}
\]
the rescaled stretching function $f_{n,t}(s)=f_n(ts)$ has
\[
          p(f_{n,t})>2^{cn},\qquad s(f_{n,t})\ge2^{cn}.
                                                               \tag{P-F62b}
\]
In particular, put
\[
 t_n=6\cdot10^{200}n^{199}+1,\qquad
 g=\gcd(n,3),\qquad
 \lambda_n=(nt_n/g),\quad \mu=(3),\quad
 \pi_n=((3t_n/g)^n).
                                                               \tag{P-F62c}
\]
For every $n\ge33$ this actual ordinary plethysm input has
positive initial coefficient, and its stretching function
$f_n(t_ns)$ is positive at every positive integer $s$.
Nevertheless, for every $0<c<1/4$, its two literal strict
indices satisfy (P-F62b) for all sufficiently large $n$.
The source parameter $h=\binom{n+2}{3}^2+n+1$ is unchanged
by the dilation. The source's total partition bit length is
\[
                  N_n=199n\log_2n+O(n),
                                                               \tag{P-F62d}
\]
so these are nonzero inputs with strict indices at least
$2^{\Omega(N_n/\log N_n)}$.
They are genuine exceptions to strict saturation in the
nonzero population of SH(b), but P-F48 proves that the entire
rescaled cubic family containing them is exponentially
sparse in that population. Thus this proposition does not
refute the density clause.
\end{proposition}

\begin{proof}
We use the exact constituent identity already proved in
P-F56, rather than infer a new pole statement for a Veronese
series. Set $M=2^{n-1}$. Let $L$ be the minimum period of
$f_n$, so that $v_2(L)=n-1$, and write its constituent
polynomials as $q_a(X)$, with subscripts taken modulo $L$.
P-F56 proves
\[
 q_a(X)-q_{a+L/2}(X)=2W_n(a),\qquad
 W_n(a)=\sum_{\substack{b\bmod M\\b\ \mathrm{odd}}}
                   A_n(\zeta^b)\zeta^{-ba},
 \qquad \zeta=e^{2\pi i/M}.
                                                               \tag{P-F62e}
\]
Every primitive-$M$ contribution here is constant as a
polynomial in $X$. All contributions with smaller two-part
have canceled in this exact identity.

The number $L$ is a period of $f_{n,t}$, with constituents
\[
                         Q_a(X)=q_{ta}(tX).
                                                               \tag{P-F62f}
\]
It is not necessary that this be its minimum period.
If a quasi-polynomial is written using another period,
each displayed constituent agrees with a minimum-period
constituent on an infinite arithmetic progression, hence
is the same polynomial. Thus the literal positivity and
saturation tests apply to all the $Q_a$ as well.
Because $t$ is odd,
$tL/2\equiv L/2\pmod L$. Substitution into (P-F62e) gives
\[
                 Q_a(X)-Q_{a+L/2}(X)=2W_n(ta).
                                                               \tag{P-F62g}
\]
This also explains why odd factors common to $t$ and $L$
cause no problem: the proof uses an existing period and
an exact identity of polynomials. It does not assume
that distinct preimage poles of a Veronese operation
cannot cancel.

For every integer $x\ge p(f_{n,t})$, all the $Q_a(x)$ are
nonnegative. The same conclusion holds whenever
$x\ge s(f_{n,t})+1$, by the definition of strict saturation,
including the identically zero constituents. Since
$Q_x(x)=f_n(tx)$, equation (P-F62g) therefore yields
\[
 f_n(tx)\ge2W_n(tx)
 \quad\text{if }x\ge p(f_{n,t})
                  \text{ or }x\ge s(f_{n,t})+1.
                                                               \tag{P-F62h}
\]
As before, one integer witness to the reverse strict
inequality proves $p(f_{n,t})>x$ and $s(f_{n,t})\ge x$.

We now produce that witness uniformly in (P-F62a).
Write $R=3g\in\{3,9\}$ and let $a_*=R^{-1}\pmod M$,
chosen between one and $M-1$. Set
\[
 Z_n=A_n(\zeta^{a_*}),\qquad
 E_n=\sum_{\substack{b\bmod M\ \mathrm{odd}\\
                              b\ne\pm a_*}}
                       |A_n(\zeta^b)|.
\]
P-F56 and P-F57 give
\[
       |Z_n|\ge2^{n^4/24-2n^3},\qquad
                 E_n\le|Z_n|/4
                      \quad\text{for all sufficiently large }n.
                                                               \tag{P-F62i}
\]
These are the original amplitudes, independent of $t$.
The coefficients of $f_n$ are real, so the two selected
terms at an argument $y$ are
$2\operatorname{Re}(Z_n\zeta^{-a_*y})$.

There is an integer $k$ with
\[
             Ra_*=1+kM,
       \qquad\frac{ta_*}{M}=\frac{tk}{R}+\frac{t}{RM}.
                                                               \tag{P-F62j}
\]
Reduction of the first equation modulo $R$ shows
$\gcd(k,R)=1$. Since $\gcd(t,3)=1$, multiplication
by $tk$ permutes the residue classes modulo $R$.
Fix any integer $X$. Among the $R$ equally spaced phases
obtained by ignoring the last summand in (P-F62j), one
has cosine at least $\cos(\pi/R)\ge1/2$, whatever the
initial phase of $Z_n\zeta^{-ta_*X}$.
For $0\le j<R$, the omitted angle has absolute value
at most $2\pi t/M$. Uniformly for $t\le2^{\epsilon n}$,
this tends to zero, because $\epsilon<1$.
For all sufficiently large $n$ it is at most $1/4$.
The inequality $|\cos(u+v)-\cos u|\le|v|$ therefore
shows that for some $x\in[X,X+R-1]$ the selected pair
contributes at least $|Z_n|/2$. The remaining terms
have total absolute value at most $E_n$. Hence
\[
       W_n(tx)\ge |Z_n|/4
       \quad\text{for some }x\in[X,X+R-1],
                                                               \tag{P-F62k}
\]
uniformly for every $X$ and every $t$ in (P-F62a).

Take $X=\lceil2^{cn}\rceil$. For the resulting witness,
$x\le2^{cn}+9$, and for all sufficiently large $n$,
\[
                  tx\le2^{(\epsilon+c)n+1}.
\]
Put $m=\binom{n+2}{3}$. The actual multiplicity is at most
the dimension of its containing symmetric power:
\[
 f_n(tx)
       \le\binom{m+(n/g)tx-1}{m-1}.
                                                               \tag{P-F62l}
\]
Using $m=n^3/6+O(n^2)$ and bounding this binomial
coefficient by $(m+ntx)^m$ gives, uniformly in $t$,
\[
 \log_2\max\{1,f_n(tx)\}
       \le\frac{\epsilon+c}{6}n^4+O_{\epsilon,c}(n^3\log n).
                                                               \tag{P-F62m}
\]
On the other hand, (P-F62i)--(P-F62k) give
\[
       \log_2(2W_n(tx))\ge n^4/24-2n^3-1.
\]
Since $(\epsilon+c)/6<1/24$, the latter is strictly
larger than (P-F62m) for all sufficiently large $n$.
Thus $f_n(tx)<2W_n(tx)$. Equation (P-F62h) proves
(P-F62b), since $x\ge2^{cn}$.

For the concrete input (P-F62c), $t_n\equiv1\pmod6$.
Its ordinary coefficient degree is
\[
        (n/g)t_n
          =\frac6g(10n)^{200}+\frac ng
                     >(10n)^{200}.
                                                               \tag{P-F62n}
\]
It also satisfies the required central divisibility
$n\mid3(n/g)t_n$. The proved quantitative support theorem
P-F54 therefore gives $f_n(t_n)>0$ for every $n\ge33$.
A nonzero homogeneous invariant in that degree has a
nonzero $s$th power in the polynomial invariant ring,
which is a domain. Thus $f_n(t_ns)>0$ for every $s\ge1$.
For any fixed $c<1/4$, choose
$0<\epsilon<1/4-c$. Since $\log_2t_n=199\log_2n+O(1)$,
we have $t_n\le2^{\epsilon n}$ for all sufficiently large
$n$, so the first part applies to this explicit input.

Finally, the source encoding sums the binary lengths
of positive partition parts. The single part of
$\lambda_n$ has length $200\log_2n+O(1)$, the inner
partition has length two, and each of the $n$ equal
parts of $\pi_n$ has length $199\log_2n+O(1)$.
This proves (P-F62d). In particular
$n=\Theta(N_n/\log N_n)$, proving the stated lower
bound in terms of actual input length. Their nonzero
initial coefficients place these examples in the
conditioned population of GCT VI Hypothesis~1.6.5(b),
while their strict saturation indices are positive.
They belong to the family $\mathcal C$ in P-F48, so
that proposition's exponentially small density bound
applies without any change of measure. No stronger
density conclusion is inferred.
\end{proof}
